\section{Implementation}
\label{sec:implementation}

\as{copied from final report -> not sure if we need to discuss these}

Here we describe some small implementation choices that had a large bearing on
performance. One was memoizing the hash function of a given term during semisort
to avoid recomputing it multiple times, which can be expensive. Hashing a term
requires hashing its function symbol, alongside the representatives of each of
its operands. This requires executing a \textsc{Find} for every term. We
therefore compute the hash of every term in the frontier at the beginning of a
round, store this in a \texttt{struct}, and perform semisort on these structs. 


Another optimization is that, when semisorting to determine congruence classes,
we perform an equality check on a secondary hash rather than deterministically
checking whether two terms are actually operand-wise equal as in the pseudocode
for \textsc{Congruent-Level1} in Figure~\ref{fig:alg:bsp}. This is because we
empirically found that having semisort probe over every bucket using this
deterministic equality check was very expensive. Changing the equality check to
a secondary hash allows the equality check when probing over a bucket to be much
cheaper. However, this introduces technical unsoundness, but this is tantamout
to a 96-bit hash collision, with probability about $10^{-14}$ for workloads of
our size (this can be computed via birthday-paradox like tricks). Although this
introduces a negligible chance of unsoundness, this change resulted in a 50\%
speedup, which we believe to be well-worth it.